\section{A locally invariant manifold near the fold and finite-interval estimates}
\label{subsec:invariant-history-details}

This section supplies the functional-analytic details used in
\cref{prop:fold-invariant-histories,prop:finite-comparison}.  We first
construct the locally invariant phase-space manifold, then solve the two
finite-interval boundary-value problems, and finally compare their
first-integral difference with the whole-line Gaussian pairings.  All norms
and constants are uniform in the number of nodes.  We use only the network
assumptions and notation of \cref{sec:rw-model-results} and the exact fold
equations in \cref{eq:fold-exact-system}.

\subsection{A smooth cutoff extension of the local fold vector field}

The graph transform requires a complete planar flow, although the resulting
histories are used only in the fold neighborhood.  We therefore extend the
local fold vector field by a smooth cutoff fixed before differentiating in
\((\nu,\eta)\).  This auxiliary extension is unrelated to the modification of
the recovery equation in \cref{subsec:global-modification}.
Introduce
\begin{equation}
 \chi=-2\alpha X,\qquad
 d=Y-\alpha X^2+\frac1{2\alpha},\qquad
 \mathcal K(X,Y)=(\chi,d).
 \label{eq:supp-fold-coordinates}
\end{equation}
In these coordinates the singular field is
\begin{equation}
 q_0^{\chi d}(\chi,d)=(1-2\alpha d,\chi d),
 \qquad \mathcal K(\gamma_0(s))=(s,0).
 \label{eq:supp-singular-field-chi-d}
\end{equation}

Put \(S=S_\delta=\sqrt{2p_0\log(1/\delta)}\), with \(p_0>4\), and let
\(R=S+B\), where \(B>2\Theta_*+2\) is fixed.  Choose
\(\omega\in C_c^\infty(\mathbb R)\), equal to one on \([-1,1]\) and
zero outside \((-2,2)\).  Fixed constants \(B_*>B+2\Theta_*+3\) and
\(d_*>0\) may be chosen so that the fold neighborhood on which the cutoff is
the identity, and all points obtained
from it by at most two delay translations lie in
\begin{equation}
 \abs{\chi}<R+B_* ,\qquad \abs d<d_*.
 \label{eq:supp-extension-region}
\end{equation}
Indeed, finite-time variational estimates for
\eqref{eq:supp-singular-field-chi-d} grow at most as \(e^{cR}\), whereas
\(\delta^\kappa e^{cR}=o(1)\) for every fixed \(0<\kappa<1\).
Define
\begin{equation}
 \Xi_\delta(\chi,d)
 =\omega\!\left(\frac{\chi}{R+B_*}\right)
  \omega\!\left(\frac d{d_*}\right),\qquad
 q_{0,\delta}=\Xi_\delta q_0^{\chi d}.
 \label{eq:supp-global-singular-field}
\end{equation}
Then \(q_{0,\delta}\) is bounded, has a complete flow, agrees with
\(q_0^{\chi d}\) on \eqref{eq:supp-extension-region}, and
\begin{equation}
 \operatorname{Lip}(q_{0,\delta})\le C(1+R).
 \label{eq:supp-singular-lipschitz}
\end{equation}

Conjugate \cref{eq:fold-exact-system} by \(\mathcal K\), replace
\(\delta\) temporarily by \(\rho\in[-\delta,\delta]\), and write the
result as
\begin{equation}
 u'=q_{0,\delta}(u)+\rho F_{N,\delta},\qquad
 \rho h'=A_Nh+\rho G_{N,\delta}.
 \label{eq:supp-extended-normal-form}
\end{equation}
Here \(F_{N,\delta}\) and \(G_{N,\delta}\) depend on the current state,
the \(m+1\) delayed states, \(\rho\), \(\nu\), and \(\eta\).  Starting
from the polynomial expressions in \cref{eq:fold-exact-system}, extend
them by fixed smooth cutoffs in the planar slots and by a fixed smooth
saturation in each transverse slot, followed by \(P_{\perp,N}\).  The
cutoffs and saturation maps are chosen to be the identity on the set in
\eqref{eq:supp-extension-region} and its twofold delay enlargement.  These
auxiliary maps are independent of \((\nu,\eta)\), and the perturbation tuple
continues to enter only through \(B_{k,N}+\Delta B_k\).

The product estimate following \cref{eq:fold-poisson-bound}, the uniform
bound on the delayed-coupling matrices, and the fixed number of delay values
give a common majorant
\begin{equation}
 A_\delta=P(R)e^{cR}
 \label{eq:supp-global-majorant}
\end{equation}
for all state and parameter derivatives used below.  In particular, for
all fixed \(a>0\) and \(j\in\mathbb N\),
\begin{equation}
 \delta^aA_\delta^j\longrightarrow0
 \qquad(\delta\to0).
 \label{eq:supp-subpower-property}
\end{equation}
No derivative of a delay evaluation is needed: the delay locations are
fixed and all parameter dependence remains in the delayed-coupling matrices.

For a complete planar field \(Q\) and a map \(H:\mathbb R^2\to E_N\),
define the phase-space embedding in fold coordinates and the finite
evaluation tuple by
\begin{align}
 \mathcal I_{Q,H}(u)(\vartheta)
 &=\bigl(\Phi_Q^\vartheta u,H(\Phi_Q^\vartheta u)\bigr),
 &&-\Theta_*\le\vartheta\le0,
 \label{eq:supp-history-lift}\\
 \mathcal E_{Q,H}(u)
 &=\bigl(\mathcal I_{Q,H}(u)(0),
          \mathcal I_{Q,H}(u)(-\theta_0),\ldots,
          \mathcal I_{Q,H}(u)(-\theta_m)\bigr).
 \label{eq:supp-history-evaluation}
\end{align}
The norm of this evaluation map is bounded by \(m+2\), independently of
\(N\).

\subsection{The invariant graph and its derivatives}

On the product of bounded continuous planar fields and bounded continuous
maps into \(E_N\), define
\begin{align}
 \mathcal T_Q(Q,H)(u)
 &=q_{0,\delta}(u)+\rho F_{N,\delta}
       (\mathcal E_{Q,H}(u);\rho,\nu,\eta),
 \label{eq:supp-graph-transform-Q}\\
 \mathcal T_H(Q,H)(u)
 &=\rho\int_0^\infty e^{A_Nr}G_{N,\delta}
       \bigl(\mathcal E_{Q,H}(\Phi_Q^{-\rho r}u);
              \rho,\nu,\eta\bigr)\,\dd r.
 \label{eq:supp-graph-transform-H}
\end{align}
The negative values of \(\rho\) are used only to take a Taylor expansion;
the RFDE graph corresponds to \(\rho=\delta>0\).

For every field in the ball
\begin{equation}
 \norm{Q-q_{0,\delta}}_\infty+\norm H_\infty
 \le C\delta A_\delta,\qquad
 \operatorname{Lip}(Q-q_{0,\delta})+\operatorname{Lip}(H)
 \le C\delta A_\delta,
 \label{eq:supp-graph-ball}
\end{equation}
the flow variational equations and \eqref{eq:supp-singular-lipschitz} give,
for each member \(\mathscr J\) of the finite derivative list used below,
\begin{equation}
 \norm{\mathscr J\Phi_Q^{\vartheta-\rho r}}
 \le A_\delta(1+r^M)e^{C(1+R)\abs{\rho}r},
 \qquad -\Theta_*\le\vartheta\le0,\quad r\ge0.
 \label{eq:supp-flow-jet-bound}
\end{equation}
The corresponding difference estimate has the same right-hand side times
the differences of the current and preceding derivative blocks.  Combining
this estimate with
\(\norm{e^{A_Nr}}\le e^{-D\gamma r}\) yields the contraction factor
\begin{equation}
 \varepsilon_\delta
 \le C\delta A_\delta
   \int_0^\infty(1+r^M)
      e^{-\{D\gamma-C\delta(1+R)\}r}\,\dd r=o(1).
 \label{eq:supp-graph-contraction}
\end{equation}
After decreasing \(\delta_0\), the ball
\eqref{eq:supp-graph-ball} is invariant and
\((\mathcal T_Q,\mathcal T_H)\) is a contraction with factor at most
\(1/2\), uniformly in \(N\), \(\nu\in I_\nu\), and \(\eta\) in a fixed ball
of admissible delayed-coupling perturbations.  We denote its fixed point by
\((Q_{N,\rho,\nu,\eta},H_{N,\rho,\nu,\eta})\).

We record why differentiating this fixed point does not lose a derivative.
For
\begin{equation}
 J_{a,b,i,j}=D_u^a\partial_\rho^b\partial_\nu^iD_\eta^j,\qquad
 \mathcal J_k=\{(a,b,i,j):b\le3,\ i,j\le2,\
                         a+b+i+j\le k\},
 \label{eq:supp-jet-list}
\end{equation}
order the blocks first by total order and then by total parameter order.
Here \(D_\eta^j\) is the symmetric Fr\'echet derivative in the norm of
\(\mathcal T_N\).  The flow variational equations and the
finite Fa\`a di Bruno formula put every current block in the form
\begin{equation}
 \mathbf J_\ell\mathcal T(Z)
 =b_\ell(\mathbf J_{<\ell}Z)
  +\rho L_\ell(\mathbf J_{<\ell}Z)[\mathbf J_\ell Z],
 \qquad \norm{\rho L_\ell}\le\varepsilon_\delta.
 \label{eq:supp-triangular-jet-equation}
\end{equation}
If a parameter derivative falls on a flow argument, the derivative of the
vector field gains one state slot and loses one parameter derivative, so it
belongs to an earlier block.  A product containing two positive-order derivatives
also contains only lower total orders.  These observations prove
\eqref{eq:supp-triangular-jet-equation}; importantly, its diagonal term
retains the factor \(\rho\).

Starting Picard iteration at \((q_{0,\delta},0)\), induction over
\(\mathcal J_{11}\) gives bounded invariant balls for all derivative
blocks.  For consecutive iterates, each block satisfies
\begin{equation}
 a_{n+1}\le\varepsilon_\delta a_n+Cq^n,
 \qquad \max\{\varepsilon_\delta,q\}<1,
 \label{eq:supp-jet-recurrence}
\end{equation}
and hence converges through \(\mathcal J_{10}\).  Uniform convergence of a
map and its derivative, followed by the fundamental theorem of calculus on
state and parameter line segments, identifies the limits as genuine
derivatives.  Thus all derivatives with
\begin{equation}
 a\le3,\qquad b\le3,\qquad i\le2,\qquad j\le2
 \label{eq:supp-required-derivatives}
\end{equation}
exist in the ranges needed here.  This argument gives full operator and
bilinear norms in \(\eta\), rather than a collection of directional
estimates.  Indeed, the bounds are first uniform on unit parameter lines;
the polarization identity turns the diagonal second derivative bound into
the symmetric bilinear bound with a dimension-independent factor.

For completeness, the finite lower-triangular system has a uniform inverse
in a scaled block norm.  Enumerate its blocks by \(0,\ldots,L\), let
\(C_{\ell k,\delta}\) bound the strict lower-block coefficient, and set
\begin{equation}
 w_{0,\delta}=1,\qquad
 w_{\ell,\delta}=1+4\sum_{k<\ell}C_{\ell k,\delta}w_{k,\delta},\qquad
 \norm V_{\rm pr,\delta}
 =\max_\ell w_{\ell,\delta}^{-1}\norm{V_\ell}.
 \label{eq:supp-prolonged-norm}
\end{equation}
Since \(\norm{(I-\rho L_\ell)^{-1}}\le2\), forward substitution gives
\begin{equation}
 \norm V_{\rm pr,\delta}\le4\norm F_{\rm pr,\delta}.
 \label{eq:supp-prolonged-inverse}
\end{equation}
For each fixed \(\delta\), the weights may depend on \(\delta\), but not on
\(N\).

Taylor's formula at \(\rho=0\), followed by the inverse coordinate change
\(\mathcal K^{-1}\), gives on the fold neighborhood where the cutoff field
equals the original equation
\begin{equation}
 Q_{N,\delta,\nu,\eta}
 =q_0+\delta q_{1,N}+\delta^2q_{2,N}
      +\delta^3\mathcal R_{3,N},
 \label{eq:supp-reduced-expansion}
\end{equation}
with the same coefficients as in \cref{eq:fold-reduced-expansion}.  The graph
component has the corresponding expansion
\begin{equation}
 H_{N,\rho,\nu,\eta}
 =\rho H_{1,N,\nu,\eta}+\rho^2\mathcal R_{H,N,\rho,\nu,\eta},
 \qquad
 H_{1,N,\nu,\eta}
 =\left.\partial_\rho H_{N,\rho,\nu,\eta}\right|_{\rho=0},
 \label{eq:supp-first-graph-coefficient}
\end{equation}
in the weighted product norm.  Split
the stable convolution at \(L\log(1/\delta)\).  The tail is bounded by
\begin{equation}
 CA_\delta\{1+L\log(1/\delta)\}^M
 e^{-D\gamma L\log(1/\delta)/2}=O(\delta^6)
 \label{eq:supp-convolution-tail}
\end{equation}
after \(L\) is fixed sufficiently large.  The finite part moves the planar
base by only \(O(\delta\log(1/\delta))\).  Applying the block induction on
a finite nested family of buffered tubes therefore gives, for
\(\ell=1,2\),
\begin{equation}
 \begin{split}
 &\norm{D_u^a\partial_\nu^iD_\eta^j
     q_{\ell,N}(\gamma_0(s)+\xi,\nu,\eta)}
 +\norm{D_u^a\partial_\nu^iD_\eta^j
     \mathcal R_{3,N}(\gamma_0(s)+\xi,\nu,\eta)}\\
 &\hspace{35mm}\le Ce^{c\abs s}\langle s\rangle^M,
 \end{split}
 \label{eq:supp-coefficient-envelope}
\end{equation}
whenever \(a\le3\), \(i,j\le2\), \(\abs s\le S+1\),
and
\begin{equation}
 \abs\xi\le C_0\delta e^{c_0\abs s}\langle s\rangle^{M_0}.
 \label{eq:supp-input-tube}
\end{equation}
The same estimate holds on the normal strip where the cutoff is the identity,
\(\mathcal K(\gamma_0(s)+\xi)=(s,d)\), \(\abs d\le\delta^\kappa\).
The same block induction applied directly to
\cref{eq:supp-graph-transform-H}, together with its prefactor \(\rho\), gives
\begin{equation}
 \norm{D_u^a\partial_\nu^iD_\eta^j
 H_{N,\delta,\nu,\eta}(\gamma_0(s)+\xi)}_N
 \le C\delta e^{c\abs s}\langle s\rangle^M
 \label{eq:supp-expanding-window-graph-bound}
\end{equation}
for \(a\le3\), \(i,j\le2\), \(\abs s\le S+1\), and \(\xi\) in
\cref{eq:supp-input-tube}.  In particular, on every fixed compact
\(s\)-interval the stable convolution gives
\begin{equation}
 \norm{H_{N,\delta,\nu,\eta}(\gamma_0(s)+\xi)}_N
 \le C\delta,\qquad \abs\xi\le C_0\delta.
 \label{eq:supp-local-graph-size}
\end{equation}

The first coefficient has zero derivative with respect to \(\eta\):
\begin{equation}
 D_\eta q_{1,N}=0.
 \label{eq:supp-q1-structural-zero}
\end{equation}
At \(\rho=0\), the only direct term involving \(\eta\) in the collective equation
is
\(K\pi_N^T\sum_k\Delta B_k\mathbf1(X_0-X_k)\), which vanishes because
\(\pi_N^T\Delta B_k=0\) for every \(k\).  The base flow and the zeroth graph are
independent of \(\eta\), proving \eqref{eq:supp-q1-structural-zero} in the
full operator norm.  The first two coefficients sample at most two delay
translations, so their values in the fold neighborhood use only the original
equations, not the global extension.

\subsection{Compatibility and injectivity of the phase-space embedding}

Fix \(\rho=\delta>0\), suppress the remaining parameters, and let
\(u(s)=\Phi_Q^su_*\), \(h(s)=H(u(s))\).  Changing variables in
\eqref{eq:supp-graph-transform-H} gives
\begin{equation}
 h(s)=\int_{-\infty}^s
 e^{A_N(s-\tau)/\delta}
 G_{N,\delta}(\mathcal E_{Q,H}(u(\tau));\delta,\nu,\eta)
 \,\dd\tau.
 \label{eq:supp-complete-convolution}
\end{equation}
Splitting this integral at \(s_0\) gives the mild transverse equation on
\([s_0,\infty)\).  Together with the first fixed-point equation, it yields
the semiconjugacy
\begin{equation}
 \mathcal S_{N,\delta,\nu,\eta}^s\mathcal I_{Q,H}(u)
 =\mathcal I_{Q,H}(\Phi_Q^su),\qquad s\ge0,
 \label{eq:supp-semiconjugacy}
\end{equation}
as long as the orbit and its delay translates remain in the region where
the extension agrees with \cref{eq:fold-exact-system}.

Evaluation at \(\vartheta=0\) is a left inverse:
\begin{equation}
 \pi_{\mathbb R^2}\mathcal I_{Q,H}(u)(0)=u.
 \label{eq:supp-present-left-inverse}
\end{equation}
Hence the phase-space embedding is injective.  To see the same fact in the original
phase space, put \(\vartheta=\delta t\).  The voltage history recovers
\begin{equation}
 X(\vartheta)=\delta^{-1}\{\pi_N^Tv(\vartheta/\delta)-1\},\qquad
 h(\vartheta)=\delta^{-2}P_{\perp,N}v(\vartheta/\delta)
                  -z_{2,N}X(\vartheta)^2,
 \label{eq:supp-recover-X-h}
\end{equation}
and the current value recovers \(Y(0)\).  The missing past values of the
current coordinate follow from the exact equation \(Y'=-X+\delta\nu\):
\begin{equation}
 Y(\vartheta)=Y(0)+\int_\vartheta^0X(r)\,\dd r
                         +\delta\nu\vartheta.
 \label{eq:supp-recover-Y-history}
\end{equation}
Thus equality of two embedded histories is equality in the full RFDE state,
not merely equality after the stationary projection.  This proves the
existence, compatibility, injectivity, expansion, and derivative assertions
used in \cref{prop:fold-invariant-histories}.

\subsection{The left and right boundary-value problems}

The finite-interval calculation needs two planar traces that agree with the
invariant graph on \([-S,S]\).  Extend its reduced field from the indicated
normal strip to a bounded \(C^3\) planar field
\(\widetilde Q_{N,\delta,\nu,\eta}\) such that it equals the graph field on
that strip and equals \(q_0\) near \(\gamma_0(\pm R)\).  A bounded linear
extension operator in the normal \(d\)-coordinate is obtained by continuing
the third-order boundary jet across each edge of the strip and multiplying
it by a cutoff supported at a fixed unscaled distance.  Fixed smooth cutoffs
then join this extension to \(q_0\) near the two endpoints.  This gives
\begin{align}
 \abs{D_z^a\partial_\nu^i
       (\widetilde Q-q_0)(\gamma_0(s)+\xi)}
 &\le C\delta e^{c\abs s}\langle s\rangle^M,
 \label{eq:supp-planar-extension-bound}\\
 \norm{D_z^aD_\eta
       (\widetilde Q-q_0)(\gamma_0(s)+\xi)}
 &\le C\delta^2e^{c\abs s}\langle s\rangle^M,
 \label{eq:supp-planar-extension-structural}
\end{align}
for \(a\le3\), \(i\le2\), and \(-R\le s\le R\).  The extension and all
cutoffs are fixed independently of \((\nu,\eta)\); using the unscaled
normal coordinate on each collar prevents inverse powers of \(\delta\).

Let \(J_a=[-R,0]\), \(J_r=[0,R]\).  For \(\sigma\in\{a,r\}\), use
the weighted norms
\begin{align}
 \norm\xi_{\mathcal X_\sigma}
 &=\sup_{s\in J_\sigma}
 e^{-c\abs s}\langle s\rangle^{-M}\abs{\xi(s)},
 \label{eq:supp-trace-state-norm}\\
 \norm f_{\mathcal Y_\sigma}
 &=\sup_{s\in J_\sigma}
 e^{-c\abs s}\langle s\rangle^{-(M-d_0)}\abs{f(s)},
 \label{eq:supp-trace-source-norm}\\
 \abs\beta_{\partial,\sigma}
 &=e^{R^2/2-cR}\langle R\rangle^{-(M-d_0)}\abs\beta,
 \label{eq:supp-trace-endpoint-norm}
\end{align}
where we fix \(d_0=3\) and take \(M\ge d_0+3\).  Put
\begin{equation}
 \mathcal E(s)=e^{s^2/2},\qquad I(s)=\int_0^s\mathcal E(t)\,\dd t,\qquad
 \tau(s)=(-1,s)^T,\qquad
 n(s)=(I(s),\mathcal E(s)-sI(s))^T.
 \label{eq:supp-fold-fundamental-pair}
\end{equation}
If \(\xi=a\tau+bn\) and \(L_0\xi=f\), then
\begin{equation}
 a'=E^{-1}\{-(E-sI)f_1+If_2\}=:A_f,\qquad
 b'=E^{-1}(sf_1+f_2)=:B_f.
 \label{eq:supp-ab-equations}
\end{equation}
The phase condition \(U(0)=0\) is \(a(0)=0\).  If the normalized normal
coordinate at the outer endpoint is \(\beta\), the two Green formulas are
\begin{align}
 a_a(s)&=-\int_s^0A_f(t)\,\dd t,
 &b_a(s)&=\beta+\int_{-R}^sB_f(t)\,\dd t,
 \label{eq:supp-left-green}\\
 a_r(s)&=\int_0^sA_f(t)\,\dd t,
 &b_r(s)&=\beta-\int_s^RB_f(t)\,\dd t.
 \label{eq:supp-right-green}
\end{align}
The estimates
\begin{equation}
 \abs{I(s)}\le C\mathcal E(s)\langle s\rangle^{-1},\qquad
 \abs{\mathcal E(s)-sI(s)}\le C\mathcal E(s)
 \label{eq:supp-I-estimates}
\end{equation}
and the one-sided Gaussian tail bound imply
\begin{equation}
 \norm\xi_{\mathcal X_\sigma}
 \le C_G\{\norm f_{\mathcal Y_\sigma}
                    +\abs\beta_{\partial,\sigma}\},
 \label{eq:supp-one-sided-green-bound}
\end{equation}
with \(C_G\) independent of \(R\).

The first integral used to normalize the endpoints is
\begin{equation}
 \mathscr H_\alpha(X,Y)
 =\frac12e^{-2\alpha Y}
   \left(\alpha Y-\alpha^2X^2+\frac12\right),\qquad
 c_{\mathscr H}=\frac2{\alpha e}.
 \label{eq:supp-first-integral}
\end{equation}
It satisfies
\begin{equation}
 c_{\mathscr H}\nabla\mathscr H_\alpha(\gamma_0(s))
 =e^{-s^2/2}(s,1)^T.
 \label{eq:supp-first-integral-adjoint}
\end{equation}
At either endpoint, the equation
\(\mathscr H_\alpha(\gamma_0+\xi)=0\) is exactly
\begin{equation}
 \mathcal E(s)b=\alpha\{-a+I(s)b\}^2.
 \label{eq:supp-endpoint-equation}
\end{equation}
It determines \(b=\mathcal B_\sigma(a)\) with
\(\mathcal B_\sigma(0)=\mathcal B_\sigma'(0)=0\), uniformly in the
endpoint norm \eqref{eq:supp-trace-endpoint-norm}.

Writing \(z=\gamma_0+\xi\), define
\begin{equation*}
 N_{N,\delta,\nu,\eta}(s,\xi)
 =\widetilde Q_{N,\delta,\nu,\eta}(\gamma_0(s)+\xi)
   -q_0(\gamma_0(s))-Dq_0(\gamma_0(s))\xi.
\end{equation*}
Let \(\mathcal G_\sigma(f,\beta)\) denote the solution given by
\cref{eq:supp-left-green,eq:supp-right-green}.  Its \(a\)-component is
independent of \(\beta\); write \(a_f(s_\sigma)\) for its value at the
outer endpoint \(s_a=-R\) or \(s_r=R\).  The trace equation, including the
nonlinear endpoint condition, is the following fixed-point equation.  Set
\(f_\xi(s)=N_{N,\delta,\nu,\eta}(s,\xi(s))\); then
\begin{equation}
 \xi=\mathcal T_\sigma(\xi)
 :=\mathcal G_\sigma\!\left(
   f_\xi,\mathcal B_\sigma(a_{f_\xi}(s_\sigma))
   \right).
 \label{eq:supp-trace-fixed-point}
\end{equation}
Equations
\eqref{eq:supp-planar-extension-bound}--\eqref{eq:supp-planar-extension-structural}
and \eqref{eq:supp-one-sided-green-bound}, together with
\(\mathcal B_\sigma'(0)=0\), show that this map on
\(\norm\xi_{\mathcal X_\sigma}\le C_0\delta\) has Lipschitz factor
\begin{equation}
 C\delta e^{cR}\langle R\rangle^M=o(1).
 \label{eq:supp-trace-contraction}
\end{equation}
It therefore has unique solutions \(z^a\) and \(z^r\) satisfying
\begin{equation}
 X^a(0)=X^r(0)=0,\qquad
 \mathscr H_\alpha(z^a(-R))
 =\mathscr H_\alpha(z^r(R))=0.
 \label{eq:supp-trace-conditions}
\end{equation}
Differentiating the same fixed-point equation gives, in full operator norm,
\begin{align}
 \abs{z^\sigma(s)-\gamma_0(s)}
 +\abs{\partial_\nu z^\sigma(s)}
 &\le C\delta W_{\rm tr}(s),
 \label{eq:supp-trace-nu-bound}\\
 \norm{D_\eta z^\sigma(s)}
 &\le C\delta^2W_{\rm tr}(s),
 \qquad \sigma\in\{a,r\},
 \label{eq:supp-trace-eta-bound}
\end{align}
where \(W_{\rm tr}(s)=e^{c_{\rm tr}\abs s}
\langle s\rangle^{M_{\rm tr}}\).  The extra power of \(\delta\) in
\eqref{eq:supp-trace-eta-bound} follows from
\eqref{eq:supp-q1-structural-zero}.  Since
\(\delta W_{\rm tr}(S)=o(\delta^\kappa)\), the fold-chart segments and all delay
translates used by the graph remain in the region where the extension agrees
with the original fold equations.  They are therefore exact RFDE histories.

\subsection{The finite-interval matching function}

Define
\begin{equation}
 D_N^{\rm fin}(\delta,\nu,\eta)
 =c_{\mathscr H}\{
     \mathscr H_\alpha(z^a(0))-
     \mathscr H_\alpha(z^r(0))\}.
 \label{eq:supp-finite-comparison-definition}
\end{equation}
For \(\ell=1,2\), let
\begin{equation}
 A_{\ell,N}(z,\nu,\eta)
 =c_{\mathscr H}\nabla\mathscr H_\alpha(z)^T
      q_{\ell,N}(z,\nu,\eta),\qquad
 A_{3,N}=c_{\mathscr H}\nabla\mathscr H_\alpha^T\mathcal R_{3,N}.
 \label{eq:supp-integrands}
\end{equation}
Put \(M_{1,N}^0=\sum_k\theta_kB_{k,N}\) and
\begin{align}
 a_{1,N}(s,\nu)
 &=e^{-s^2/2}\left\{
    \nu+\frac{\kappa_N}{8\alpha^3}s^4+d_N^{\rm del}s\right\},
 &d_N^{\rm del}&=-\frac K{2\alpha}\pi_N^TM_{1,N}^0\mathbf1,
 \label{eq:supp-leading-integrand}\\
 a_{2,N}(s)[\eta]
 &=-e^{-s^2/2}s^2\Lambda_N(\eta).
 \label{eq:supp-structural-integrand}
\end{align}
Equations \eqref{eq:fold-parameter-pairing},
\eqref{eq:fold-return-source}, and
\eqref{eq:supp-first-integral-adjoint} give, on \(\abs s\le S\),
\begin{equation}
 A_{1,N}(\gamma_0(s),\nu,\eta)=a_{1,N}(s,\nu),\qquad
 D_\eta A_{2,N}(\gamma_0(s),\nu,0)[\eta]=a_{2,N}(s)[\eta].
 \label{eq:supp-central-integrand-identities}
\end{equation}
The whole-line pairings are
\begin{equation}
 \int_{\mathbb R}a_{1,N}(s,\nu)\,\dd s
 =\sqrt{2\pi}(\nu-\nu_{0,N}),\qquad
 \int_{\mathbb R}a_{2,N}(s)[\eta]\,\dd s
 =-\sqrt{2\pi}\Lambda_N(\eta).
 \label{eq:supp-whole-line-pairings}
\end{equation}

Let \(z^{\rm pw}=z^a\) on \([-R,0]\) and \(z^{\rm pw}=z^r\) on
\([0,R]\), keeping the two values at zero separate when integrating.
Also set
\begin{equation}
 B_{2,N}^S(\nu)=
 \int_{-S}^SA_{2,N}(\gamma_0(s),\nu,0)\,\dd s.
 \label{eq:supp-baseline-second-order}
\end{equation}

\begin{proposition}[Finite-interval Gaussian comparison]
\label[proposition]{prop:finite-gap-bridge}
On a common neighborhood of \(I_\nu\times\{0\}\) in
\(I_\nu\times\mathcal T_N\), the matching function has the
exact decomposition
\begin{align}
 D_N^{\rm fin}
 ={}&\delta\sqrt{2\pi}(\nu-\nu_{0,N})
    +\delta^2B_{2,N}^S(\nu)
    -\delta^2\sqrt{2\pi}\Lambda_N(\eta)\notag\\
 &+E_{\partial,N}+E_{{\rm tail},N}+E_{{\rm mov},N}
   +E_{{\rm nl},N}+E_{3,N}+E_{{\rm strip},N},
 \label{eq:supp-exact-comparison-decomposition}
\end{align}
where the six error terms are defined in
\eqref{eq:supp-error-boundary}--\eqref{eq:supp-error-strip} below.  Uniformly
in \(N\),
\begin{align}
 E_{\partial,N}&=0,
 \label{eq:supp-comparison-boundary-bound}\\
 \abs{E_{{\rm tail},N}}+\abs{\partial_\nu E_{{\rm tail},N}}
 +\norm{D_\eta E_{{\rm tail},N}}
 &\le C\delta^4,
 \notag\\
 \abs{E_{{\rm strip},N}}+\abs{\partial_\nu E_{{\rm strip},N}}
 +\norm{D_\eta E_{{\rm strip},N}}&\le C\delta^4,
 \label{eq:supp-comparison-tail-bound}\\
 \abs{E_{{\rm mov},N}}+\abs{\partial_\nu E_{{\rm mov},N}}
 &\le C\delta^2,
 &\norm{D_\eta E_{{\rm mov},N}}&\le C\delta^3,
 \label{eq:supp-comparison-moving-bound}\\
 \abs{E_{3,N}}+\abs{\partial_\nu E_{3,N}}
 +\norm{D_\eta E_{3,N}}&\le C\delta^3,
 \label{eq:supp-comparison-third-order-bound}\\
 \abs{E_{{\rm nl},N}}+\abs{\partial_\nu E_{{\rm nl},N}}
 &\le C\delta^2\norm{\eta}_{\mathcal T_N}^2,
 &\norm{D_\eta E_{{\rm nl},N}}&
 \le C\delta^2\norm{\eta}_{\mathcal T_N}.
 \label{eq:supp-comparison-nonlinear-bound}
\end{align}
Moreover, \(B_{2,N}^S\) and its first \(\nu\)-derivative are uniformly
bounded.  Consequently,
\begin{align}
 \frac{D_N^{\rm fin}}\delta
 &=\sqrt{2\pi}(\nu-\nu_{0,N})+O(\delta),
 \label{eq:supp-finite-level}\\
 \partial_\nu D_N^{\rm fin}
 &=\delta\sqrt{2\pi}+O(\delta^2),
 \label{eq:supp-finite-nu-slope}\\
 \norm{D_\eta D_N^{\rm fin}
          +\delta^2\sqrt{2\pi}\Lambda_N}
 &\le C\{\delta^3+\delta^2\norm{\eta}_{\mathcal T_N}\}.
 \label{eq:supp-finite-structural-slope}
\end{align}
\end{proposition}

\begin{proof}
Since \(\mathscr H_\alpha\) is a first integral of \(q_0\), the
fundamental theorem of calculus and \eqref{eq:supp-trace-conditions} give
\begin{align}
 D_N^{\rm fin}
 ={}&E_{\partial,N}
 +c_{\mathscr H}\int_{-R}^0
  \nabla\mathscr H_\alpha(z^a(s))^T
       \widetilde Q(z^a(s))\,\dd s\notag\\
 &+c_{\mathscr H}\int_0^R
  \nabla\mathscr H_\alpha(z^r(s))^T
       \widetilde Q(z^r(s))\,\dd s,
 \label{eq:supp-fundamental-comparison-identity}
\end{align}
where
\begin{equation}
 E_{\partial,N}=c_{\mathscr H}\{
 \mathscr H_\alpha(z^a(-R))-\mathscr H_\alpha(z^r(R))\}=0.
 \label{eq:supp-error-boundary}
\end{equation}
There are no parameter-dependent integration limits.  In particular,
\eqref{eq:supp-error-boundary} and all its \((\nu,\eta)\)-derivatives
vanish exactly.

On \([-S,S]\), insert \eqref{eq:supp-reduced-expansion} into
\eqref{eq:supp-fundamental-comparison-identity}, add and subtract the
singular-orbit integrands, and use Taylor's formula in \(\eta\).  This yields
\eqref{eq:supp-exact-comparison-decomposition} with
\begin{align}
 E_{{\rm tail},N}
 ={}&-\delta\int_{\abs s>S}a_{1,N}(s,\nu)\,\dd s
     -\delta^2\int_{\abs s>S}a_{2,N}(s)[\eta]\,\dd s,
 \label{eq:supp-error-tail}\\
 E_{{\rm mov},N}
 ={}&\delta\int_{-S}^S{}^{\rm pw}
  \{A_{1,N}(z^{\rm pw}(s))-A_{1,N}(\gamma_0(s))\}\,\dd s\notag\\
 &+\delta^2\int_{-S}^S{}^{\rm pw}
  \{A_{2,N}(z^{\rm pw}(s),\eta)
             -A_{2,N}(\gamma_0(s),\eta)\}\,\dd s,
 \label{eq:supp-error-moving}\\
 E_{{\rm nl},N}
 ={}&\delta^2\int_{-S}^S
 \{A_{2,N}(\gamma_0(s),\eta)-A_{2,N}(\gamma_0(s),0)
       -D_\eta A_{2,N}(\gamma_0(s),0)[\eta]\}\,\dd s,
 \label{eq:supp-error-nonlinear}\\
 E_{3,N}
 ={}&\delta^3\int_{-S}^S{}^{\rm pw}
       A_{3,N}(z^{\rm pw}(s),\delta,\nu,\eta)\,\dd s,
 \label{eq:supp-error-third-order}\\
 E_{{\rm strip},N}
 ={}&c_{\mathscr H}
 \int_{[-R,-S]\cup[S,R]}{}^{\rm pw}
 \nabla\mathscr H_\alpha(z^{\rm pw}(s))^T
 \{\widetilde Q(z^{\rm pw}(s))-q_0(z^{\rm pw}(s))\}\,\dd s.
 \label{eq:supp-error-strip}
\end{align}
The unused arguments \((\nu,\eta)\) in
\eqref{eq:supp-error-moving}--\eqref{eq:supp-error-third-order} have been
suppressed.

It remains to bound these terms.  Direct differentiation of
\eqref{eq:supp-first-integral} and
\eqref{eq:supp-trace-nu-bound} gives, for each fixed \(k\le4\),
\begin{equation}
 \norm{D^k\mathscr H_\alpha(\gamma_0(s)+\xi)}
 \le Ce^{-s^2/2}e^{c\abs s}\langle s\rangle^M
 \label{eq:supp-first-integral-envelope}
\end{equation}
on the trace tube.  The product of this Gaussian with any finite collection
of the weights in \eqref{eq:supp-coefficient-envelope},
\eqref{eq:supp-trace-nu-bound}, and
\eqref{eq:supp-trace-eta-bound} is integrable on \(\mathbb R\).

The mean-value theorem now gives the first two bounds in
\eqref{eq:supp-comparison-moving-bound}.  For the derivative with respect to
\(\eta\) of
the first integral in \eqref{eq:supp-error-moving}, the direct derivative of
\(A_{1,N}\) vanishes by \eqref{eq:supp-q1-structural-zero}, while
\(D_\eta z^{\rm pw}=O(\delta^2W_{\rm tr})\).  The second integral has the
additional prefactor \(\delta^2\), and the difference of its direct
\(\eta\)-derivatives at \(z^{\rm pw}\) and \(\gamma_0\) contains
\(z^{\rm pw}-\gamma_0=O(\delta W_{\rm tr})\).  Hence
\(\norm{D_\eta E_{{\rm mov},N}}\le C\delta^3\).
The same chain-rule estimate with
\eqref{eq:supp-coefficient-envelope} proves
\eqref{eq:supp-comparison-third-order-bound}.

Taylor's formula with integral remainder in the Banach space
\(\mathcal T_N\) gives
\begin{equation}
 A_{2,N}(\eta)-A_{2,N}(0)-D_\eta A_{2,N}(0)[\eta]
 =\int_0^1(1-t)D_{\eta\eta}^2A_{2,N}(t\eta)[\eta,\eta]\,\dd t.
 \label{eq:supp-structural-taylor-formula}
\end{equation}
The Gaussian envelope for the second Fr\'echet derivative proves
\eqref{eq:supp-comparison-nonlinear-bound}, including its first
\(\nu\)- and \(\eta\)-derivatives.

Finally, for fixed \(c,M\),
\begin{equation}
 \int_{\abs s>S}e^{-s^2/2+c\abs s}\langle s\rangle^M\,\dd s
 \le Ce^{-S^2/2+cS}\langle S\rangle^{M+1}
 =\delta^{p_0-o(1)}.
 \label{eq:supp-gaussian-tail}
\end{equation}
The explicit formulas \eqref{eq:supp-leading-integrand} and
\eqref{eq:supp-structural-integrand} therefore give the stated bound for
\(E_{{\rm tail},N}\).  On the two endpoint strips, which have fixed length,
\eqref{eq:supp-planar-extension-bound},
\eqref{eq:supp-planar-extension-structural}, and
\eqref{eq:supp-first-integral-envelope} give the same Gaussian bound for
\(E_{{\rm strip},N}\).  Since \(p_0>4\), both are bounded by
\(C\delta^4\) after decreasing \(\delta_0\).

The coefficient envelope makes \(B_{2,N}^S\) and its first
\(\nu\)-derivative uniformly bounded.  Substitution of all preceding
estimates into \eqref{eq:supp-exact-comparison-decomposition} proves
\eqref{eq:supp-finite-level}--\eqref{eq:supp-finite-structural-slope},
which are the estimates asserted in \cref{prop:finite-comparison}.
\end{proof}
